\documentclass[8pt]{article}
\usepackage[a4paper, total={6in, 8in}]{geometry}
\usepackage{graphicx} % Required for inserting images
\usepackage{amsthm}
\usepackage{amsmath}
\usepackage{amsfonts}
\usepackage{amsopn}
\usepackage{comment}
\usepackage{amssymb}
\usepackage{hyperref}
\usepackage{bussproofs}
\usepackage[all, 2cell]{xy}
\usepackage[all]{xy}
\usepackage{rotating}
\usepackage{lscape}
\usepackage{minted}
\usepackage{dsfont}
\usepackage{tikz}
\usepackage{tikz-cd}
\usepackage{stmaryrd}
\usepackage{multicol}
\usepackage{cmll}
\usepackage{wasysym}
\usetikzlibrary{shapes.geometric, arrows, positioning, backgrounds}
\usetikzlibrary{cd}
\makeatletter
\tikzcdset{
  eq node/.style={
    commutative diagrams/math mode=false, anchor=center},
  eq/.style={
    phantom,
    /tikz/every to/.append style={
      edge node={node[commutative diagrams/eq node]
        {\@eqnswtrue\make@display@tag\ltx@label{#1}}}}}}
\makeatother

\theoremstyle{definition}
\newtheorem{definition}{Definition}[section]

\theoremstyle{definition}
\newtheorem{theorem}{Theorem}[section]

\theoremstyle{definition}
\newtheorem{claim}{Claim}[section]

\theoremstyle{definition}
\newtheorem{ex}{Example}[section] 

\theoremstyle{definition}
\newtheorem{cons}{Construction}[section] 

\theoremstyle{definition}
\newtheorem{rem}{Remark}[section] 


\theoremstyle{definition}
\newtheorem{prop}{Proposition}[section]

\theoremstyle{definition}
\newtheorem{lemma}{Lemma}[section]

\theoremstyle{definition}
\newtheorem{fact}{Fact}[section]

\theoremstyle{definition}
\newtheorem{remark}{Remark}[section]

\theoremstyle{definition}
\newtheorem{notation}{Notation}[section]

\theoremstyle{definition}
\newtheorem{example}{Example}[section]

\theoremstyle{definition}
\newtheorem{exercise}{Exercise}[section]

\theoremstyle{definition}
\newtheorem{col}{Corollary}[section]

\theoremstyle{question}
\newtheorem{question}{Question}

\let\strokeL\L
\renewcommand\L{\mathbf{L}}

\DeclareMathOperator{\bang}{!}

\DeclareMathOperator*{\colim}{\operatorname{colim}}
\newcommand{\Ob}[1]{\operatorname{Ob}({\mathcal{#1}})}
\newcommand{\Cov}[1]{\operatorname{Cov}(#1)}

\newcommand{\circvee}{\mathbin{\text{\ooalign{\hfil$\vee$\hfil\cr\kern0.1ex$\bigcirc$}}}}
\newcommand{\circwedge}{\mathbin{\text{\ooalign{\hfil$\wedge$\hfil\cr\kern0.1ex$\bigcirc$}}}}

\title{$h$-adic topology and profinite linear lambda terms}
\author{Daniel Rogozin}
\date{ }

\newcommand{\IntHom}[1]{[#1]}

\begin{document}

\maketitle
\tableofcontents

\section{Profinite Lambda terms and related concepts}

One starts with the observation that if we have a work over a finite alphabet $\Sigma$:
\[
W = a_{w_1} \ldots a_{w_n} \in \Sigma^*
\]
then one can represent $W$ as a well-typed term in the simply typed lambda calculus:
\[
a_{w_1} : q \to q, \ldots, a_{w_n} : q \to q \vdash W : q \to q
\]
where $q$ is a base type that we informally interpret as a type of states, so we have a closed lambda
term
\[
\lambda a_{w_1}^{q \to q} \ldots \lambda a_{w_n}^{q \to q} \lambda x^q. a_{w_k} (\ldots (a_{w_1} \: q)\ldots) : (q \to q) \ldots (q \to q) \to q \to q
\]
encoding a word $w$ via the instance of the Church encoding. Further, we will call the type of such terms $\operatorname{Church}_{\Sigma}$.
Further, we have the monoid isomorphism:
\[
\Sigma^* \cong \{ \text{$M$ is a lambda term} \: | M : \operatorname{Church}_{\Sigma} \:\}/\beta \eta
\]
In other words, words over a finite alphabet can be viewed as lambda terms of a particular type factorised module the $\beta\eta$-equivalence.
Generally, there is amount literature studying both syntactic and semantic aspects of the Church encoding of the finite words.

The claim made in \cite{van2023profinite} connects the finitary interpretation of STLC in the category of finite sets ${\bf FinSet}$
with deterministic finite automata as follows. Take any finite set $Q$,
so for any type $A$, we define a set $[\![A]\!]_Q$ inductively defined as
\begin{itemize}
\item $[\![p]\!]_Q = Q$
\item $[\![A \to B]\!]_Q = [\![A]\!]_Q \Rightarrow [\![B]\!]_Q$
\end{itemize}
so every term $x_1 : A_1, \ldots, x_n : A_n \vdash M : A$ transports
to the function between finite sets $[\![M]\!] : [\![A_1]\!]_Q \times \ldots \times [\![A_n]\!]_Q \to [\![A]\!]_Q$.
On closed terms, the interpretation induces a function $[\![.]\!]_Q  : \Lambda_{\beta \eta}(A) \to [\![A]\!]_Q$
mapping every closed term $M$ of type $A$ to its interpretation $[\![M]\!]_Q$. 
Further, a word $W = a_{w_1} \ldots a_{w_k} \in \Sigma^*$ is mapped to the following morphism in ${\bf FinSet}$:
$[\![W]\!]_Q : (Q \Rightarrow Q) \times \ldots \times (Q \Rightarrow Q) \to (Q \Rightarrow Q)$, i.e.,
\[
[\![W]\!]_Q : f_{w_1}, \ldots, f_{w_k} \mapsto f_{w_k} \circ \ldots \circ f_{w_1}
\]

Recall that a \emph{deterministic finite state automaton} on a finite alphabet $\Sigma = \{ a_1, \ldots, a_n \}$
is a structure $\mathcal{A} = (Q, \delta, q_0, Acc)$, where $Q$ is a finite set of states,
$\delta : \Sigma \times Q \to Q$ is a transition function, $q_0 \in Q$ is the initial state and
$Acc \subseteq Q$ is the set of accepting states. The transition function $\delta$ induces a family
of functions $\delta_{a_i}(q_k) = \delta(a_i, q_k) : Q \to Q$ for each $a_i \in \Sigma$ and $q_k \in Q$.
$\delta_{a_i}(q_k) = q_{k'}$ means that $\mathcal{A}$ transitions from $q_k$ to $q_k'$ encountering $a_i$.
Now one can apply the interpretation $[\![W]\!]_Q$ to the functions $\delta_{a_{w_i}}$ for $1 \leq i \leq k$,
so we obtain
\[
\delta_w : [\![W]\!]_Q (\delta_{a_{w_1}}, \ldots, \delta_{a_{w_k}}) : Q \to Q
\]

Therefore, given a deterministic finite automaton $\mathcal{A} = (Q, \delta, q_0, Acc)$,
the tuple $(Q, \delta, q_0)$ induces the evaluation function
\[
\operatorname{eval}_{\delta, q_0} : [\![\operatorname{Church}_{\Sigma}]\!]_Q \to Q
\]
transporting every functional $F \in [\![\operatorname{Church}_{\Sigma}]\!]_Q$ to the state
$F(\delta_{a_{w_1}}, \ldots, \delta_{a_{w_k}})$ in $Q$, therefore we have
\[
\Sigma^{*} \cong \Lambda_{\beta\eta}(\operatorname{Church}_{\Sigma}) \to [\![\operatorname{Church}_{\Sigma}]\!]_Q \xrightarrow{\operatorname{eval}_{\delta, q_0}} Q
\]
so we associate every finite word $w \in \Sigma^*$ with the final state $q_f = \delta_w (q_0)$.
The inverse image $Acc \subseteq Q$ under this composite is called a \emph{regular language} $L_{\mathcal{A}}$, the
set of finite words recognisable by $\mathcal{A}$.

Regular languages $L_{\mathcal{A}}$ form a Boolean algebra 
$\operatorname{Reg}_Q(\operatorname{Church}_{\Sigma})$ of regular languages of lambda terms
of type $\operatorname{Church}_{\Sigma}$ recognised by $Q$. The algebra can be defined as the image 
of the Boolean algebra homomorphism 
$[\![.]\!]^{-1}_Q : 2^{[\![\operatorname{Church}_{\Sigma}]\!]} \to 2^{\Sigma^*}$ obtained by applying
the contravariant powerset functor $2^{(.)} : {\bf FinSet}^{op} \to {\bf BA}$ to 
the semantic bracket $[\![.]\!]_Q$. The Boolean algebra $\operatorname{Reg}_Q(A)$ of regular languages of
lambda terms of any higher order type $A$ recognisable by a finite set $Q$ is defined as the image
of the Boolean algebra homomorphism $[\![.]\!]^{-1}_Q : 2^{[\![A]\!]_Q} \to 2^{\Lambda_{\beta \eta}(A)}$.
That is, a set $L$ of lambda terms of type $A$ is recognisable by the finite set $Q$ exactly when it has
the form $[\![.]\!]^{-1}_Q(\operatorname{Acc}) = \{ M \in \Lambda_{\beta\eta}(A) \: | \: [\![M]\!]_Q \in \operatorname{Acc} \}$
for some choice $\operatorname{Acc} \subseteq [\![A]\!]_Q$ of accepting elements. Thus, the collection
$\operatorname{Reg}(A) \subseteq 2^{\Lambda_{\beta \eta}(A)}$ of lambda terms of type $A$ as follows
\[
\operatorname{Reg}(A) = \bigcup \{ \operatorname{Reg}_Q(A) \: | \: Q \in {\bf FinSet} \}.
\]

The monoid $\widehat{\Sigma^*}$ is the free profinite monoid generated by $\Sigma$. $\widehat{\Sigma^*}$ can be constructed
as the limit, computed in the category of monoids ${\bf Mon}$, of the projective system of finite monoid
homomorphisms
\[
(\Sigma^*/\phi \to \Sigma^*/\phi')_{\phi \subseteq \phi'}
\]
where $\phi$ and $\phi'$ are finite index congruences on the free monoid $\Sigma^*$ such that $\phi \subseteq \phi'$.
Every finite index congruence $\phi$, in turn, induces a surjective homomorphism $h : \Sigma^* \to M$
for the quotient monoid $M = \Sigma^*/\phi$. The surjectivity condition on $h$ can be relaxed to show that
$\widehat{\Sigma^*}$ of profinite words is the codirected limit of the composite functor
\[
\Sigma^*/{\bf FinMon} \xrightarrow{\pi} {\bf FinMon} \hookrightarrow {\bf Mon}
\]
$\Sigma^*/{\bf FinMon}$ is the slice category whose objects $(M, h)$ are pairs consisting
of a finite monoid $M$ and a homomorphism $h : \Sigma^* \to M$ and whose morphisms 
$(M, h) \to (M',h')$ are homomorphisms $f : M \to M'$ such that the following triangle commutes:
\[
\begin{tikzcd}
  & \Sigma^* \ar[dl, "h"'] \ar[dr, "h'"] \\
M \ar[rr, "f"'] && M'
\end{tikzcd}
\]

The projection functor $\pi : \Sigma^*/{\bf FinMon} \rightarrow {\bf FinMon}$ transports $(M, h)$
to its underlying monoid $M$. One obtains $\widehat{\Sigma^*}$ as the limit of a codirected diagram 
of monoid homomorphisms $(M \to M')_{(M,h) \to (M',h')}$ extending the diagram formed from 
the family $(\Sigma^*/\phi \to \Sigma^*/\phi')_{\phi \subseteq \phi'}$ to all homomorphisms 
$h : \Sigma^* \to M$ to a finite monoid $M$.

Recall that every homomorphism $h : \Sigma^* \to M$ induces a deterministic finite automaton 
with $Q := M$, $\delta(a, q) = q \cdot_M a$ for every $q \in M$ and $a \in \Sigma$. 
So every homomorphism $h : \Sigma^* \to M$ (assume $M = \{ q_1, \ldots, q_m\}$) induces a decomposition of $\Sigma^*$
into $m$ components $L_{q_i} = h^{-1}(q_i)$ for $1 \leq i \leq m$ where each $L_{q_i}$ is a regular language.
Let $\operatorname{Reg}_{(M,h)}(\Sigma)$ denote the Boolean algebra of languages of the form $L_q = h^{-1}(q)$
for $q \in M$. We have a functor:
\[
\operatorname{Reg}_{(.)}(\Sigma) : (\Sigma^*/{\bf FinMon})^{op} \to {\bf BA}
\]
The Boolean algebra $\operatorname{Reg}_{(M,h)}(\Sigma)$ coincides with $2^h$, the image of $h : \Sigma^* \to M$ 
applied to the contravariant powerset functor. Also, the following directed diagram in ${\bf BA}$ associated to
$\operatorname{Reg}_{(.)}(\Sigma)$:
\[
(\operatorname{Reg}_{(M',h')}(\Sigma) \to \operatorname{Reg}_{(M,h)}(\Sigma))_{(M,h) \to (M',h')}
\]
can be obtained by applying the contravariant powerset functor to the above trianlge and $(M \to M')_{(M,h) \to (M',h')}$.

\subsection{The Explicit Description of Profinite Topology}

This description is based on \cite[C.X.2]{pin2010mathematical}. Let $\Sigma$ be
a finite alphabet and let $M$ be a finite monoid, given $u, v \in \Sigma^*$,
a monoid homomorphism $f : \Sigma^* \to M$ \emph{separates} $u$ and $v$ if $f(u) \neq f(v)$. A monoid
$M$ separates $u$ and $v$ if there is a homomorphism $\Sigma^* \to M$ separating $u$ and $v$. The important 
fact is that every $u$ and $v$ are separated by some monoid $M$. 

A metric on $\Sigma^*$ is defined as follows: two words are closed if they are separated by the large monoid.
Take $u, v \in \Sigma^*$, put:
\begin{itemize}
  \item $r(u,v) = \min \{ |M| \: | \: \text{$M$ separates $u$ and $v$}\}$,
  \item $d(u, v) = 2^{-r(u,v)}$.
\end{itemize}
Therefore, $d(u,v) < 2^{n}$ iff $u$ and $v$ cannot be separated by a monoid smaller than $n$
and $d(u,v) < \varepsilon$ iff $u$ and $v$ cannot be separated by a monoid of size $n$,
where $n$ is the unique natural number such that $2^{-(n+1)} < \varepsilon \leq 2^{-n}$.
$d$ is a metric on $\Sigma^{*}$ such that $(\Sigma^*, d)$ is an ultrametric space such that, for each 
$u, v, w \in \Sigma^*$:
\begin{itemize}
  \item $d(uw, vw) \leq d(uv)$,
  \item $d(wu, wv) \leq d(uv)$.
\end{itemize}
As a topological space, $(\Sigma^*, d)$ is discrete, but open balls are of more interest. For any $n \geq 2$,
consider the congruence relation $\sim_n$ on $\Sigma^*$ defined as:
\[
u \sim_n v \Leftrightarrow \forall M \: (|M| \leq n \rightarrow (()\forall f : \Sigma^* \twoheadrightarrow M) \rightarrow f(u) = f(v)))
\]
Clearly $\sim_n$ is a congruence of finite index as they are finitely many surjections $\Sigma^* \twoheadrightarrow M$
for a fixed finite $M$. Thus $d(u,v) < 2^{-n}$ iff they cannot be separated by a monoid of size smaller than $n$.

\begin{prop} Let $\Sigma$ be a finite alphabet,
\begin{enumerate}
  \item The concatenation function $(u,v) \mapsto uv$ is uniformly continuous,
  \item Every open ball of  $(\Sigma^*, d)$  is a recognisable language and every recognisable language is a finite union
  of open balls.
  \item $(\Sigma^*, d)$ is totally bounded: for every $n > 0$, $\Sigma^*$ is covered by a finite number
  of open balls of radius strictly smaller than $2^{-n}$.
\end{enumerate}
\end{prop}

\begin{prop}
  Let $\Sigma_1$ and $\Sigma_2$ be finite alphabets, then $f : \Sigma_1 \to \Sigma_2$
  is uniformly continuous iff for every recognisable language of $\Sigma_2$ $L$, $f^{-1}(L)$
  is also recognisable.
\end{prop}

Let $\widehat{\Sigma^{*}}$ denote the completion of $\Sigma^*$, whose points are called profinite words of $\Sigma$.

\begin{prop}
  Let $M$ be a finite monoid and let $f : \Sigma^* \to M$ be uniformly continuous, then $f$ uniquely extends
  to $\widehat{f} : \widehat{\Sigma^*} \to M$.
\end{prop}
The following is the universal property of the free profinite monoid:
\begin{prop} Let $\Sigma$ be a finite alphabet, then
  if $f : \Sigma \to M$ is a function to a finite monoid $M$, there exists a unique uniformly
  continuous monoid morphism $\widehat{f} : \widehat{\Sigma^*} \to M$
  such that for each $a \in \Sigma$, $\widehat{f}(a) = f(a)$. $\widehat{f}$ is onto iff $f(\Sigma)$
  generates $M$.
\end{prop}
\begin{prop}
  Let $\Sigma_1$ and $\Sigma_2$ be finite alphabets, then every uniformly continuous morphism
  $f : \Sigma_1^* \to \Sigma_2^*$ uniquely extends to uniformly continuous 
  $\widehat{f} : \widehat{\Sigma_1^*} \to \widehat{\Sigma_2^*}$.
\end{prop}

The metric on $\widehat{\Sigma^*}$ is defined similarly. Let $M$ be a finite monoid, then
a morphism $f : \Sigma^* \twoheadrightarrow M$ separates profinite words $u, v \in \widehat{\Sigma^*}$
if $\widehat{f}(u) = \widehat{f}(v)$. For $u, v \in u, v \in \widehat{\Sigma^*}$, we set:
\begin{itemize}
  \item $r(u,v) = \min \{ |M| \: | \: \text{$M$ separates $u$ and $v$} \}$,
  \item $d(u, v) = 2^{-r(u,v)}$.
\end{itemize}
\begin{prop}
  The profinite topology on $\widehat{\Sigma^*}$ is the least topology 
  making every continuous every morphism $\Sigma^*$ onto a finite discrete monoid.
\end{prop}

\subsection{Profinite Lambda Terms}

One moves from profinite words to profinite lambda terms in the following way. For any finite set $Q$, one has the interpretation
of any type $A$ as the finite set $[\![A]\!]_Q$. Given another finite set $Q'$ with a relation $R \subseteq Q \times Q'$, 
one can relate different interpretations of $A$ by constructing a relation $[\![A]\!]_R \subseteq [\![A]\!]_Q \times [\![A]\!]_{Q'}$.
Such relations are called \emph{logical relations}. Lambda terms are parametric: that is, if $M$ has type $A$ and $R \subseteq Q \times Q'$,
then $([\![M]\!]_Q, [\![M]\!]_{Q'}) \in [\![A]\!]_R$.

Now let us discuss how one can define the exponential of relations. Let $S \subseteq P \times P'$ and
$R \subseteq Q \times Q'$ be relations on finite sets, the exponential $S \Rightarrow R \subseteq (P \Rightarrow Q) \times (P' 
\Rightarrow Q')$ is given as:
\[
S \Rightarrow R := \{ (g,h) \: | \: \forall x \in P \forall y \in P' \: (x S y \to g(x) R h(y))\}
\]
Therefore, for any relation $R \subseteq Q \times Q'$ between finite sets $Q$ and $Q'$, one constructs a relation
$[\![A]\!]_R \subseteq [\![A]\!]_Q \times [\![A]\!]_{Q'}$ defined by induction on $A$ as follows:
\begin{itemize}
  \item $[\![p]\!] = R$,
  \item $[\![A \Rightarrow B]\!]_R = [\![A]\!]_R \Rightarrow [\![B]\!]_R$
\end{itemize}
The \emph{fundamental lemma of logical relations} states that for any $M \in \Lambda_{\beta\eta}(A)$ for any $A$
and for any $R \subseteq Q \times Q'$, the interpretations of $M$ are related as follows: 
\[
([\![M]\!]_Q, [\![M]\!]_{Q'}) \in [\![A]\!]_R
\]
Further, we will be using partial surjections since they are stable under taking relational exponentials. Also,
if $f : Q \twoheadrightarrow Q'$ is a partial surjection, so $[\![A]\!]_f \subseteq [\![A]\!]_Q \times [\![A]\!]_{Q'}$
for any $A$. Besides, we have the inclusion of Boolean algebras $\operatorname{Reg}_{Q'}(A) \subseteq \operatorname{Reg}_Q(A)$,
which is proved using the fundamental lemma of logical relations.

\begin{theorem}
  The diagram of Boolean algebras $\operatorname{Reg}_{(.)}(A) : {\bf FinPSurj}^{op} \to {\bf BA}$:
  \[
  (\operatorname{Reg}_{Q'}(A) \hookrightarrow \operatorname{Reg}_{Q}(A))_{f : Q \to Q' \in {\bf FinPSurj}}
  \]
  is directed, and its colimit in ${\bf BA}$ coincides with the Boolean algebra $\operatorname{Reg}(A)$
  of regular languages for a higher-order type $A$. 
\end{theorem}
Also, for any simple type $A$, the set of regular languages $\operatorname{Reg}(A)$ is closed under finite intersections.

Now let us discuss how the space of profinite lambda terms is constructed. The idea is that for any
simple type $A$ as parametric families of semantic elements, so they form a Stone space dual to 
the Boolean algebra $\operatorname{Reg}(A)$.

\begin{prop}
  Let $Q, Q'$ be finite sets,

  \begin{enumerate} 
    \item The set of atoms $\operatorname{At}_Q(A)$ of the finite Boolean algebra $\operatorname{Reg}_Q(A)$
    is in bijection with the set $[\![A]\!]_{Q}^{\bullet}$ of definable elements of simple type $A$, given by the
    function $q \mapsto [\![.]\!]^{-1}_Q(\{ q \}) : \operatorname{At}_Q(A) \to [\![A]\!]_{Q}^{\bullet}$.
    \item For partial surjection $f : Q \twoheadrightarrow Q'$ and for any simple type $A$,
    the set $[\![A]\!]_Q^{\bullet} \subseteq [\![A]\!]_Q$ of definable elements is contained 
    in the domain of $[\![A]\!]_f$, and the restriction of $[\![A]\!]_f$ to $[\![A]\!]_Q^{\bullet}$
    is the unique function $p_{Q,Q'}$ such that the following triangle commutes:
    \[
    \begin{tikzcd}
      & \Lambda_{\beta\eta}(A) \ar[dl, "{[\![.]\!]}_Q^{\bullet}"'] \ar[dr, "{[\![.]\!]}_{Q'}^{\bullet}"] \\
    {[\![A]\!]}_Q^{\bullet} \ar[rr, "p_{Q,Q'}"'] && {[\![A]\!]}_{Q'}^{\bullet}
    \end{tikzcd}
    \]
  \end{enumerate}
\end{prop}

Note that if $f, g : Q \to Q'$ are partial surjections, then, while the semantic interpretations
$[\![A]\!]_f, [\![A]\!]_g : [\![A]\!]_Q \to [\![A]\!]_Q'$ are different, their restrictions 
to $[\![A]\!]^{\bullet}_Q$ are both equal to $p_{Q,Q'}$.

One defines profinite lambda terms of simple type $A$ as follows.
\begin{definition}
  Let $A$ be a simple type, the set of \emph{profinite lambda terms} is defined as the limit 
  $\widehat{\Lambda}_{\beta\eta}(A)$ in ${\bf Set}$ of the diagram:
  \[
  ([\![A]\!]_f^{\bullet} : [\![A]\!]_Q^{\bullet} \twoheadrightarrow [\![A]\!]_{Q'}^{\bullet})_{f : Q \twoheadrightarrow Q' \in {\bf FinPSurj}}
  \]
\end{definition}



\section{$h$-adic topology}

The following summary is based on \cite[Chapter XVI]{kassel2012quantum}.

Consider the complex algebra $K = \mathbb{C}[[h]]$ of complex formal series
in one variable $h$. Every element $f \in \mathcal{K}$ has the form
\[
f = \sum \limits_{n = 0}^{\infty} a_n h^n
\]
where $a_n \in \mathbb{C}$ for each $n \in \mathbb{N}$. Every polynomial $f \in \mathbb{C}[h]$ is clearly a formal series. $K$ form a commutative ring with operations given as
\begin{itemize}
  \item $f + g = \sum \limits_{n = 0}^{\infty} a_n h^n + \sum \limits_{n = 0}^{\infty} b_n h^n = \sum \limits_{n = 0}^{\infty} (a_n + b_n) h^n$
  \item $f \cdot g = (\sum \limits_{n = 0}^{\infty} a_n h^n) \cdot (\sum \limits_{n = 0}^{\infty} b_n h^n) = \sum \limits_{n=0}^{\infty} (\sum \limits_{n = p + q} a_p b_q) h^n$
\end{itemize}
A formal series $f$ is invertible iff its every coefficient is non-zero.

For any $n > 0$, consider the algebra $K_n = \mathbb{C}[h]/(h^n)$
of truncated polynomials obtained as the quotient of the ring of complex
polynomials generated by the ideal $(h^n)$. There is a morphism
\[
\pi_n : K \to K_n
\]
sending each formal series $\sum \limits_{n = 0}^{\infty} a_n h^n$
to the class of $\sum \limits_{k = 0}^{n - 1} a_k h^k$ modulo $(h^n)$.
$\pi_n$ is onto and its kernel is the ideal $h^n K$ generated by $h^n$
in the ring of formal series. Therefore, we have the following isomorphism:
\[
\mathbb{C}[[h]]/(h^n) \cong \mathbb{C}[h]/(h^n)
\]
For $n > 0$, there is a surjective map of algebras $p_n : K_n \to K_{n - 1}$
generated by the inclusion of ideals $(h^{n}) \subset (h^{n-1})$. 
So we have the inverse system of algebras $(K_n, p_n)_{n > 0}$ and its inverse
limit $\varprojlim \limits_{n} K_n$. One has $p_n \circ \pi_n = \pi_{n - 1}$
for each $n$. Also, there is a unique morphism of algebras from $K$ to
$\varprojlim \limits_{n} K_n$ whose composition
with the projection of the inverse limit is equal to $\pi_n$. 
Besides, the map $\pi : \mathbb{C}[[h]] \twoheadrightarrow \varprojlim \limits_{n} \mathbb{C}[[h]] / (h^n)$
is the isomorphism, which is partly the instance of Birkhoff’s subdirect representation theorem, \cite[Theorem 8.6]{sankappanavar1981course}.

Let us equip $K$ with the \emph{$h$-adic topology}. The family of 
of subsets $\pi^{-1}_n(U_n)$ for $n > 0$ and $U_n \subseteq K_n$ form an open basis on $K$.
$\{ 0 \}$ is a family of open neighbourhoods in the discrete set $K_n$, 
the family $\pi_n^{-1}(0)$ is a family of open neighbourhoods of $0$ in $K$ for the $h$-adic topology.
There is also the structure of a metric space on $K$ given as follows. Given a formal series
$f = \sum \limits_{n = 0}^{\infty} a_n h^n$, let $o(f)$ denote the unique natural number 
such that $a_{o(f)} \neq 0$ and $a_k = 0$ for each $k < o(f)$. For $f = 0$, set $o(f) = +\infty$
and extend $(\mathbb{N}, <)$ by adding the top element $+\infty$. It is easy to see that:
\[
(h^n) = \{ f \in K \: | \: o(f) > n - 1\}
\]
and therefore
\[
\bigcap \limits_{n > 0} (h^n) = \{ 0 \}
\]
We also have for any series $f, g \in K$:
\[
o(f + g) \geq \min (o(f), o(g))
\]
Now define $|.| : K \to \mathbb{R}_{\geq 0}$ as
\[
|f| = 2^{-o(f)}
\]
Then if we put $d(f,g) = |f - g|$, then $(\mathbb{C}[[h]], d)$ becomes an ultrametric space.

Let $M$ be a left $K$-module, consider the family $(h^n M)_{n > 0}$ of submodules of $M$
with the canonical $K$-linear projections:
\[
p_n : M_n = M/h^n M \to M_{n - 1} = M/h^{n - 1} M
\]
They form an inverse system of $K$-modules, so we can consider the inverse limit
\[
\widetilde{M} = \varprojlim \limits_{n} M_n
\]
having a natural structure of a left $K$-module. $\widetilde{M}$ has a natural topology as well,
the inverse limit topology, for which the family of submodules $(h^n \widetilde{M})_n$ form an open basis.
The projections $i_n : M \to M_n$ induce a unique $K$-linear map $i : M \to \widetilde{M}$
such that $\pi_n \circ i = i_n$ for each $n$. The kernel of $i$ is given as
\[
\operatorname{Ker}(i) = \bigcap \limits_{n > 0} h^n M
\]

A $K$-module $M$ is \emph{separated} if $\operatorname{Ker}(i) = \{ 0 \}$. It is \emph{complete} 
if $i$ is injective.

For any module $K$, the module $M/\operatorname{Ker}(i)$ is separated and the completion 
$\widetilde{M}$ of $M$ is complete. 

Generally, any separated complete $K$-module $M$ is equipped with the $h$-adic topology induced
by the inverse limit topology on $\widetilde{M}$ via $M \cong \widetilde{M}$.

Now let us describe the following class of separated complete $K$-modules. Let $V$ be a vector space over $\mathbb{C}$,
let $V[[h]]$ denote the set of all formal series in variable $h$
\[
f = \sum \limits_{n = 0}^{\infty} a_n h^n
\]
where $a$'s range over elements of $V$. There is a left $K$-module on $V[[h]]$, and such modules are called 
\emph{topologically free modules}. One can observe that any topologically free module is separated and complete. Also,
if $\{ e_i \}_{i \in I}$ is a basis of $V$, then the $K$-submodule generated by $\{ e_i \}_{i \in I}$ is dense 
in $V$ for the $h$-adic topology. Besides, let $N$ be a complete separated $K$-module, there is a natural isomorphism
in the category of $K$-modules and $K$-linear maps:
\[
\operatorname{Hom}_K(V[[h]], N) \cong \operatorname{Hom}_K(V, N)
\]
Further, a $K$-module $M$ is topologically free iff it is complete, separated and torsion-free, that is,
there is such that $m h \neq 0$ for any non-zero $m \in M$.

Let $M$ and $N$ be left modules over the algebra $K = \mathbb{C}[[h]]$. The $K$-module 
$M \otimes_K N$ is the quotient of the vector space $M \otimes N$ obtained by spanning all
elements of the form $fm \otimes n - m \otimes fn$ for $f \in K$, $m \in M$ and $n \in N$.

\begin{definition}
  The topological tensor product $M \widetilde{\otimes} N$ of $M$ and $N$ is the $h$-adic completion of $M \otimes_K N$:
  \[
  M \widetilde{\otimes} N := \widetilde{M \otimes_K N} = \varprojlim \limits_{n > 0} (M \otimes_K N) / h^n (M \otimes_K N)
  \]
\end{definition}
By definition, $M \widetilde{\otimes} N$ is complete.

Take $m \in M$ and $n \in N$, let $m \widetilde{\otimes} n$ denote the image of $m \otimes n$ under the following image:
\[
M \otimes N \rightarrow M \otimes_K N \rightarrow M \widetilde{\otimes} N
\]
The subspace of the topological tensor product spanned by all elements of this form is dense in $M \otimes N$.

\begin{prop}
  \item The class left modules over the algebra $K = \mathbb{C}[[h]]$ with $\widetilde{\otimes}$ is endowed with the structure of a symmetric monoidal category
  \item If $M$ and $N$ be topologically free modules, so is $M \widetilde{\otimes} N$, that is:
  \[
  V[[h]] \widehat{\otimes} W[[h]] = (V \otimes W)[[h]].
  \]
\end{prop}

\section{On Dialogue Categories and their internal language}
\begin{definition}
  A dialogue category is a monoidal category $(\mathcal{C}, \otimes, e)$ equipped with an object $\bot$ along with a family
  of isomorphisms
  \[
  \varphi_{A,B} : \mathcal{C}(A \otimes B, \bot) \cong \mathcal{C}(B, A \multimap \bot)
  \]
  natural in $x$ and $y$; along with a family of isomorphisms
  \[
  \psi_{A, B} : \mathcal{C}(A \otimes B, \bot) \cong \mathcal{C}(A, \bot \multimapinv B)
  \]
  natural in $x$ and $y$ as well.
\end{definition}

\subsection{Term Assignment for Tensorial Logic}

\begin{figure}[H]
  \hrule
  \caption{Term Calculus for Tensorial Logic}~\label{fig:1}


  \begin{minipage}[b]{0.33333\textwidth}
    \begin{small}
      \begin{prooftree}
        \AxiomC{$ $}
        \RightLabel{$\mathds{1}{\bf I}$}
        \UnaryInfC{$\vdash {\bf 1} : \mathds{1}$}
      \end{prooftree}
    \end{small}
    \end{minipage}%
    \begin{minipage}[b]{0.33333\textwidth}
      \begin{small}
        \begin{prooftree}
          \AxiomC{$ $}
          \RightLabel{\bf ax}
          \UnaryInfC{$x : A \vdash x : A$}
        \end{prooftree}
        \end{small}
    \end{minipage}%
    \begin{minipage}[b]{0.33333\textwidth}
    \begin{small}
      \begin{prooftree}
        \AxiomC{$\Gamma \vdash M : \mathds{1}$}
        \AxiomC{$\Delta \vdash N : B$}
        \RightLabel{$\mathds{1}{\bf E}$}
        \BinaryInfC{$\Gamma, \Delta \vdash {\bf let} \: {\bf 1} = M \: {\bf in} \: N : B$}
      \end{prooftree}
    \end{small}
    \end{minipage}
  \hfill
  \begin{minipage}{0.45\textwidth}
    \begin{small}
    \begin{prooftree}
      \AxiomC{$x : A, \Gamma \vdash M : \bot $}
      \RightLabel{${\multimap \bot}_{\bf I}$}
      \UnaryInfC{$\Gamma \vdash \nabla_l x. M : A \multimap \bot $}
    \end{prooftree}

      \begin{prooftree}
      \AxiomC{$\Gamma, x : A \vdash M : \bot $}
      \RightLabel{${\bot \multimapinv}_{\bf I}$}
      \UnaryInfC{$\Gamma \vdash \nabla_r x. M : \bot \multimapinv A$}
    \end{prooftree}

    \begin{prooftree}
      \AxiomC{$\Gamma \vdash M : A$}
      \AxiomC{$\Delta \vdash N : B$}
      \RightLabel{$\otimes{\bf I}$}
      \BinaryInfC{$\Gamma, \Delta \vdash M \otimes N : A \otimes B$}
    \end{prooftree}
  \end{small}
  \end{minipage}
  \begin{minipage}{0.45\textwidth}
    \begin{tabular}{p{\textwidth}}
      \begin{small}
      \begin{prooftree}
        \AxiomC{$\Delta \vdash N : A$}
        \AxiomC{$\Gamma \vdash M : A \multimap \bot$}
        \RightLabel{$\multimap{\bf E}$}
        \BinaryInfC{$\Gamma, \Delta \vdash N \triangleleft M : \bot$}
      \end{prooftree}

      \begin{prooftree}
        \AxiomC{$\Gamma \vdash M : A \multimapinv \bot$}
        \AxiomC{$\Delta \vdash N : A$}
        \RightLabel{$\multimap{\bf E}$}
        \BinaryInfC{$\Gamma, \Delta \vdash M \triangleright N : \bot$}
      \end{prooftree}


      \begin{prooftree}
        \AxiomC{$\Gamma \vdash M : A \otimes B$}
        \AxiomC{$\Delta, x : A, y : B \vdash N : C$}
        \RightLabel{$\otimes{\bf E}$}
        \BinaryInfC{$\Delta, \Gamma \vdash {\bf let} \: x \otimes y \: = \:M \: {\bf in} \: N : C$}
      \end{prooftree}
    \end{small}
    \end{tabular}
  \end{minipage}
  \vspace{\baselineskip}
  \hrule
\end{figure}

\begin{definition} Reduction rules for tensorial logic:
  \begin{itemize}
    \item[$(\multimap \bot \beta)$] $N \triangleleft (\nabla_l x. M) \to M[x := N]$,
    \item[$(\bot \multimapinv \beta)$] $(\nabla_r x. M) \triangleright N \to M [x:=N]$,
    \item[$(\mathds{1}\beta)$] ${\bf let} \: {\bf 1} = {\bf 1} \: {\bf in} \: N \to N$,
    \item[$(\otimes\beta)$] ${\bf let} \: x \otimes y = M \otimes N \: {\bf in} \: P \to P [x := M, y := N]$,
    \item[$(\multimap \bot \eta)$] $(\nabla_l x. x \triangleleft M) \to M$ if $x$ is not free in $M$,
    \item[$(\bot \multimapinv \eta)$] $(\nabla_r x. M \triangleright x) \to M$ if $x$ is not free in $M$,
    \item[$(\otimes \eta)$] ${\bf let} \: x \otimes y = M \: {\bf in} \: N[z := x \otimes y] \to N [z:= M]$,
    \item[$(\otimes_{\bf nat})$] $P[w := {\bf let} \: x \otimes y = M \: {\bf in} \: N] \to {\bf let} \: x \otimes y = M \: {\bf in} \: P [w := N]$,
    \item[$(\mathds{1}_{\bf nat})$] $P [w := {\bf let} \: {\bf 1} = M \: {\bf in} \: N] \to {\bf let} \: {\bf 1} = M \: {\bf in} \: P[w := N]$.
  \end{itemize}
\end{definition}

\subsection{Parametricity for Tensorial Logic}

The following subsection develops parametricity for the tensorial calculus in the fashion of \cite{aberle2025substructural}
and \cite{zhao2010relational}.

\bibliographystyle{alpha} 
\bibliography{Text}

\end{document}